\question What would Python display? Run these statements in order.
If a statement causes an error, identify the error and continue.

\begin{blocksection}
\begin{lstlisting}
>>> t = Tree(5, [Tree(3, [Tree(8)]), Tree(1), Tree(7)])
>>> t.label
\end{lstlisting}
\begin{solution}[.35in]
\begin{lstlisting}
5
\end{lstlisting}
\end{solution}
\end{blocksection}

\begin{blocksection}
\begin{lstlisting}
>>> t.branches[0].label
\end{lstlisting}
\begin{solution}[.35in]
\begin{lstlisting}
3
\end{lstlisting}
\end{solution}
\end{blocksection}

\begin{blocksection}
\begin{lstlisting}
>>> is_leaf(t.branches[1])
\end{lstlisting}
\begin{solution}[.35in]
\begin{lstlisting}
True
\end{lstlisting}
\end{solution}
\end{blocksection}

\begin{blocksection}
\begin{lstlisting}
>>> len(t.branches)
\end{lstlisting}
\begin{solution}[.35in]
\begin{lstlisting}
3
\end{lstlisting}
\end{solution}
\end{blocksection}

\begin{blocksection}
\begin{lstlisting}
>>> t.branches[1].branches[0]
\end{lstlisting}
\begin{solution}[.35in]
\begin{lstlisting}
IndexError: list index out of range
\end{lstlisting}
\end{solution}
\end{blocksection}

\begin{blocksection}
\begin{lstlisting}
>>> t.branches[0].branches
\end{lstlisting}
\begin{solution}[.35in]
\begin{lstlisting}
[Tree(label=8, branches=[])]
\end{lstlisting}
\end{solution}
\end{blocksection}

\begin{blocksection}
\begin{lstlisting}
>>> u = Tree(2, [t.branches[0]])
>>> u.branches[0].branches[0].label
\end{lstlisting}
\begin{solution}[.35in]
\begin{lstlisting}
8
\end{lstlisting}
\end{solution}
\end{blocksection}

\begin{blocksection}
\begin{lstlisting}
>>> print(u)
\end{lstlisting}
\begin{solution}[1in]
\begin{lstlisting}
2
  3
    8
\end{lstlisting}
\end{solution}
\end{blocksection}

\begin{questionmeta}
\textbf{Teaching Tips}\par\nopagebreak[4]
Suggested Time: 6--8 min; Difficulty: Easy
\nopagebreak[4]
\begin{itemize}
    \item \textbf{Name each value:} \lstinline{t.branches} is a list; \lstinline{t.branches[i]} is a \lstinline{Tree}. Have students state the type of each intermediate expression before taking another attribute.
    \item \textbf{Leaves have no branches:} A leaf's \lstinline{branches} is the empty list, so indexing into it raises an \lstinline{IndexError}. Contrast with linked lists, where the empty case \lstinline{()} is not a \lstinline{Link} at all.
    \item \textbf{Reading print:} Depth is shown by indentation; students can redraw the tree from the printout.
\end{itemize}
\end{questionmeta}
